\section{Discussion}
\label{sec:discussion}

This work extends the HydroGym paper's GPPO-vs-PPO result on Kolmogorov
flow~\cite{hydrogym2026} in two directions, and the second produces a
seed-dominance result that is itself the main finding.

At the regular operating point ($\Re=45$) the paper's core claim holds: GPPO reaches
PPO's best reward in substantially fewer episodes in 4 of 5 seeds
($\approx$\SIrange{63}{84}{\percent}, $\text{seed}=0$'s \SI{63}{\percent}
closest to the reported \SI{65}{\percent}), the direction is seed-stable
across all five, and the off-policy family reaches better final reward than
the on-policy family in every seed. GSAC never beats SAC here --- the
analytic gradient does not help the off-policy method when the flow is not
chaotic. This is consistent with the SAPO~\cite{sapo2025} reading that the
analytic gradient earns its keep by \emph{stabilizing} a hard landscape
rather than by accelerating an easy one; at $\Re=45$ SAC's replay buffer
already solves the task.

At the chaotic operating point ($\Re=300$), by contrast, no separable
ordering survives: the conference version, at $n=2$, reported that GSAC is
the most robust of the four and that analytic-gradient injection delays the
break in direction.
Over five seeds the second claim does not hold at all (GPPO$-$PPO mean
$-2.4$ episodes, sign-flip $p=0.81$) and the first survives only as a
non-significant lean (every pairwise permutation $p\ge 0.19$;
Kruskal--Wallis $p=0.57$, Friedman $p=0.77$). Each algorithm's first-NaN
episode spans more than an order of magnitude; PPO, first to break at
$\text{seed}=0$, is the most robust of the four at $\text{seed}=2$; at
$\text{seed}=2$ both off-policy methods collapse at episode 2. The ``clean
multiplicative'' $n=1$ story and the ``GSAC most robust'' $n=2$ story are
both small-$n$ artifacts. The finding that replaces them: in the
chaotic regime the breaking point is dominated by the random path of
training, not by the algorithm --- within-algorithm cross-seed s.d.\
(30--57 episodes) swamps the spread of between-algorithm means (35--75).
Analytic-gradient injection, off-policy structure, and their combination
fail to impose a statistically separable ordering on when a differentiable
chaotic-flow simulation diverges under RL training. This is a concrete,
quantified instance of the chaos-induced gradient instability of Metz et
al.~\cite{metz2021} and the Re-shrunk predictability horizon of
Boffetta--Musacchio~\cite{boffetta2017}: at this problem scale the
instability is not merely large, it is \emph{irreducibly stochastic}. It
also matches what RPO~\cite{rpo2025} targets --- the training instability of
BPTT-through-a-surrogate --- and suggests that the RPO/GI-PPO line of
adaptive weighting, rather than a fixed mix, is the right tool for this
regime.

Two mechanistic results do survive, both algorithm-design-independent. The
first is buffer dilution
(\cref{sec:buffer}): a post-trigger survival window has non-zero probability
only when the poisoned fraction of the replay buffer at the trigger is
$\sim$1\% or below --- that is, only for GSAC runs that trigger past
$\sim$episode 120 --- and exactly one of the 20 $\Re=300$ runs shows such a
window, in precisely that case, so the mechanism is suggestive rather than
confirmed. The second is the early-init failure mode (\cref{sec:earlyinit}):
when the randomly initialized policy composed with the $\Re=300$ base flow
already sits near the solver's stability boundary, divergence occurs within
about two episodes, before any actor-learning algorithm has a foothold; bare
SAC and GSAC fail identically here, with normal control amplitude and a
normal analytic-gradient norm at the point of divergence. Together these
explain why GSAC has the highest ceiling of the four (141, 63, and 122
episodes when it is given a runway) but no reliable floor (episode 2 at
$\text{seed}=2$).

For a practitioner adding an analytic through-simulator gradient to an RL
controller for a near-turbulent flow: (i) expect no reliable improvement in
training stability from the gradient alone, and budget for divergence at an
unpredictable episode; (ii) an off-policy buffer buys a recovery chance only
if the divergence happens late, when the buffer is large --- it is not a
general safety net; and (iii) the useful engineering artifacts are the
diagnostics (first-NaN-episode distribution, pre-break flatness check,
buffer-dilution estimate), not a ranking of algorithms. None of this refutes
the HydroGym paper's core claim that GPPO helps at its tested operating
point; it shows the benefit does not carry a seed-robust ordering once the
regime is chaotic.

Several caveats bound these conclusions. The study uses five seeds: this is
enough to falsify the $n=2$ orderings and to support the negative claim, but
GSAC's non-significant median lean ($p\approx0.19$--$0.38$) would need a
substantially larger seed budget to confirm or reject --- though, because the
finding is one of seed-dominance, additional seeds are expected to strengthen
it rather than overturn it. The chaotic regime is $\Re=300$ rather than the
originally intended $\Re=450$, forced by the solver's $\Re\approx350$--$380$
stability limit at this grid and time step; a milder chaotic regime might
admit a more stable ordering, which the planned Reynolds sweep
$\{45,100,200,300\}$ is designed to test. The GPPO/GSAC mixing ratio,
horizon, and remaining hyperparameters are our own choices, since the
HydroGym paper does not publish them, so the \SI{63}{\percent} versus
\SI{65}{\percent} agreement is a strong corroborating signal but not a claim
of identical implementation. The PPO update is deliberately minimal (one
full-batch step per episode), which makes ``PPO'' here weaker than a tuned
implementation; the comparison is therefore relative and budget-matched
across all four algorithms rather than absolute. The episodes (15 RL steps)
and the analytic horizon (6 steps) are both short and may be insufficient to
observe the full course of RL convergence or of chaotic divergence, and the
50/50 analytic/model-free mix is held fixed --- the adaptive $\alpha$-graded
estimator~\cite{suh2022}, AHAC's~\cite{ahac2024} adaptive horizon, and
SAPO's~\cite{sapo2025} entropy-based smoothing are not tested. Finally, the
two mechanistic results rest on thin evidence: the buffer-dilution
calculation of \cref{sec:buffer} is supported by a single recovery window in
20 runs, and the early-init failure of \cref{sec:earlyinit} by a single
seed; both are real but rare, and neither is established as a rate.
